%% kapa: exam=1 round=28 year=2017 period=2 form=A subject=accounting subject_ko=회계학 seq=16 number=56
%% source: Q-Net 감정평가사 기출문제 (cst00309) artlSeq=5199603 file=extracted/1008/5199603/2017년 제28회 감정평가사 1차시험 문제지/2교시 A형.pdf
%% conversion: kordoc 4.24.51 -> tools/gen1.py (auto); review=proofread-sonnet; flags=none
\begin{question}{56}
% [원문 그대로] 영위하고 하고 있는 (중복 표기는 원문)
자동차부품 제조업을 영위하고 하고 있는 (주)감평은 20x1년 초 임대수익 목적으로 건물(취득원가 ￦1,000,000, 잔여 내용연수 5년, 잔존가치 ￦0, 정액법 감가상각)을 취득하였다. 한편, 20x1년 말 동 건물의 공정가치는 ￦1,200,000이다. 다음 설명 중 옳지 않은 것은? (단, 해당 건물은 매각예정으로 분류되어 있지 않다.)
\begin{choices}
  \item 원가모형을 적용할 경우, 20x1년 감가상각비는 ￦200,000이다.
  \item 공정가치모형을 적용할 경우, 20x1년 감가상각비는 ￦200,000이다.
  \item 공정가치모형을 적용할 경우, 20x1년 평가이익은 ￦200,000이다.
  \item 공정가치모형을 적용할 경우, 20x1년 당기순이익은 ￦200,000만큼 증가한다.
  \item 공정가치모형을 적용할 경우, 20x1년 기타포괄손익에 미치는 영향은 ￦0이다.
\end{choices}
\end{question}
\answer{2}
